\documentclass[unicode, 12pt, a4paper,oneside,fleqn]{article}

\usepackage{../src/styles}
\usepackage{pgfplots}
\pgfplotsset{compat=1.18}
\usepackage{alltt}

\begin{document}

\begin{titlepage}
\begin{center}
\bfseries

{\Large Московский авиационный институт\\ (национальный исследовательский университет)
}

\vspace{48pt}

{\large Факультет информационных технологий и прикладной математики
}

\vspace{36pt}

{\large Кафедра вычислительной математики и~программирования
}

\vspace{48pt}

Лабораторная работа \textnumero 2 по курсу \enquote{Дискретный анализ}

\end{center}

\vspace{72pt}

\begin{flushright}
\begin{tabular}{rl}
Студент: & И.\,А. Казаков \\
Преподаватель: & А.\,А. Кухтичев \\
Группа: & М8О-206БВ-24 \\
Дата: & \\
Оценка: & \\
Подпись: & \\
\end{tabular}
\end{flushright}

\vfill

\begin{center}
\bfseries
Москва, \the\year
\end{center}
\end{titlepage}

\pagebreak

\CWHeader{Лабораторная работа \textnumero 2}

\CWProblem{
    Необходимо создать программную библиотеку, реализующую структуру данных \textbf{PATRICIA tree}. На её основе разработать программу-словарь. Ключи — регистронезависимые строки (до 256 символов), значения — числа от 0 до $2^{64}-1$. 
    
    Программа должна поддерживать:
    \begin{itemize}
        \item Добавление, поиск и удаление элементов.
        \item Сохранение и загрузку дерева в компактном бинарном виде.
    \end{itemize}
    
    Требуется провести исследование качества программы с использованием профилировщика \texttt{gprof} и средств анализа памяти.
}

\pagebreak

\section{Описание структуры данных}
Дерево \textbf{PATRICIA} (Practical Algorithm to Retrieve Information Coded in Alphanumeric) представляет собой сжатое префиксное дерево. В отличие от обычного дерева, здесь узлы с одним потомком удаляются, а ветвление происходит по номеру критического бита (\texttt{bit\_idx}), в котором ключи различаются. Это позволяет хранить ключи произвольной длины, не увеличивая высоту дерева пропорционально длине строк.

\section{Исходный код}
Реализация выполнена на языке C++.

\begin{lstlisting}[language=C++]
#include <algorithm>
#include <cstdint>
#include <cstring>
#include <fstream>
#include <iostream>
#include <string>
#include <type_traits>

template <typename T>
concept Numeric = std::is_scalar_v<T>;

template <Numeric T>
class TrieNode {
   public:
    std::string k;
    T val;
    size_t bit_idx;
    TrieNode *l_ptr, *r_ptr;

    TrieNode(const std::string& key, T value, size_t idx)
        : k(key), val(value), bit_idx(idx), l_ptr(nullptr), r_ptr(nullptr) {}

    ~TrieNode() = default;
};

template <Numeric T>
class TrieMap {
   private:
    TrieNode<T>* head = nullptr;

    uint8_t fetch_bit(const std::string& key, size_t bit_pos) {
        if (bit_pos == 0) {
            return 0;
        }
        size_t char_idx = (bit_pos - 1) / 8;
        if (char_idx >= key.length()) {
            return 0;
        }
        return (static_cast<uint8_t>(key[char_idx]) >> (7 - (bit_pos - 1) % 8)) & 1;
    }

    size_t diff_at(const std::string& s1, const std::string& s2) {
        size_t max_len = std::max(s1.length(), s2.length());
        for (size_t i = 0; i < max_len; ++i) {
            uint8_t b1 = (i < s1.length()) ? (uint8_t)s1[i] : 0;
            uint8_t b2 = (i < s2.length()) ? (uint8_t)s2[i] : 0;
            if (b1 != b2) {
                uint8_t xor_res = b1 ^ b2;
                int p = 7;
                while (!(xor_res & (1 << p))) {
                    p--;
                }
                return i * 8 + (7 - p) + 1;
            }
        }
        return 0;
    }

    void deep_clear(TrieNode<T>* node) {
        if (node->l_ptr && node->l_ptr->bit_idx > node->bit_idx) {
            deep_clear(node->l_ptr);
        }
        if (node->r_ptr && node->r_ptr->bit_idx > node->bit_idx) {
            deep_clear(node->r_ptr);
        }
        delete node;
    }

    void write_recursive(std::ofstream& os, TrieNode<T>* node) {
        if (!node) {
            return;
        }
        size_t sz = node->k.size();
        os.write(reinterpret_cast<const char*>(&sz), sizeof(sz));
        os.write(node->k.c_str(), sz);
        os.write(reinterpret_cast<const char*>(&node->val), sizeof(T));

        if (node->l_ptr && node->l_ptr->bit_idx > node->bit_idx) {
            write_recursive(os, node->l_ptr);
        }
        if (node->r_ptr && node->r_ptr->bit_idx > node->bit_idx) {
            write_recursive(os, node->r_ptr);
        }
    }

    void locate(const std::string& key, TrieNode<T>*& target, TrieNode<T>*& parent, TrieNode<T>*& grandparent) {
        target = head->l_ptr;
        parent = head;
        grandparent = nullptr;
        while (target->bit_idx > parent->bit_idx) {
            grandparent = parent;
            parent = target;
            if (fetch_bit(key, target->bit_idx) == 0) {
                target = target->l_ptr;
            } else {
                target = target->r_ptr;
            }
        }
    }

   public:
    TrieMap() = default;

    ~TrieMap() {
        if (!head) {
            return;
        }
        if (head->l_ptr != head) {
            deep_clear(head->l_ptr);
        }
        delete head;
    }

    bool add(const std::string& key, T value) {
        if (!head) {
            head = new TrieNode<T>(key, value, 0);
            head->l_ptr = head;
            return true;
        }

        TrieNode<T>* prev = head;
        TrieNode<T>* curr = head->l_ptr;
        while (curr->bit_idx > prev->bit_idx) {
            prev = curr;
            if (fetch_bit(key, curr->bit_idx) == 0) {
                curr = curr->l_ptr;
            } else {
                curr = curr->r_ptr;
            }
        }

        if (key == curr->k) {
            return false;
        }

        size_t new_idx = diff_at(key, curr->k);
        prev = head;
        curr = head->l_ptr;
        while (curr->bit_idx > prev->bit_idx && curr->bit_idx < new_idx) {
            prev = curr;
            if (fetch_bit(key, curr->bit_idx) == 0) {
                curr = curr->l_ptr;
            } else {
                curr = curr->r_ptr;
            }
        }

        TrieNode<T>* leaf = new TrieNode<T>(key, value, new_idx);
        uint8_t bit = fetch_bit(key, new_idx);
        if (bit == 0) {
            leaf->l_ptr = leaf;
            leaf->r_ptr = curr;
        } else {
            leaf->l_ptr = curr;
            leaf->r_ptr = leaf;
        }

        if (prev == head) {
            prev->l_ptr = leaf;
        } else {
            if (fetch_bit(key, prev->bit_idx) == 0) {
                prev->l_ptr = leaf;
            } else {
                prev->r_ptr = leaf;
            }
        }
        return true;
    }

    T* find(const std::string& key) {
        if (!head) {
            return nullptr;
        }
        TrieNode<T>* p = head;
        TrieNode<T>* c = head->l_ptr;
        while (c->bit_idx > p->bit_idx) {
            p = c;
            if (fetch_bit(key, c->bit_idx) == 0) {
                c = c->l_ptr;
            } else {
                c = c->r_ptr;
            }
        }
        if (c->k == key) {
            return &(c->val);
        } else {
            return nullptr;
        }
    }

    bool del_key(const std::string& key) {
        if (!head) {
            return false;
        }
        if (key == head->k && head->l_ptr == head) {
            delete head;
            head = nullptr;
            return true;
        }

        TrieNode<T>*target, *q, *parent_q;
        locate(key, target, q, parent_q);

        if (!target || target->k != key) {
            return false;
        }

        if (target == q && target != head) {
            TrieNode<T>* child;
            if (target->l_ptr == target) {
                child = target->r_ptr;
            } else {
                child = target->l_ptr;
            }

            if (fetch_bit(key, parent_q->bit_idx) == 0) {
                parent_q->l_ptr = child;
            } else {
                parent_q->r_ptr = child;
            }
            delete target;
            return true;
        }

        TrieNode<T>*t_dummy, *q_back, *q_back_p;
        locate(q->k, t_dummy, q_back, q_back_p);

        TrieNode<T>* q_child;
        if (q->l_ptr == target) {
            q_child = q->r_ptr;
        } else {
            q_child = q->l_ptr;
        }

        if (q_child == q) {
            q_child = target;
        }

        if (parent_q) {
            if (fetch_bit(key, parent_q->bit_idx) == 0) {
                parent_q->l_ptr = q_child;
            } else {
                parent_q->r_ptr = q_child;
            }
        } else {
            head->l_ptr = q_child;
        }

        if (q_back && q_back != q) {
            if (fetch_bit(q->k, q_back->bit_idx) == 0) {
                q_back->l_ptr = target;
            } else {
                q_back->r_ptr = target;
            }
        }

        target->val = q->val;
        target->k = q->k;
        delete q;
        return true;
    }

    int store(const std::string& path) {
        std::ofstream os(path, std::ios::binary);
        if (!os) {
            std::cerr << "ERROR: " << std::strerror(errno) << "\n";
            return -1;
        }
        if (head) {
            write_recursive(os, head->l_ptr);
        }
        return 0;
    }

    int restore(const std::string& path) {
        std::ifstream is(path, std::ios::binary);
        if (!is) {
            std::cerr << "ERROR: " << std::strerror(errno) << "\n";
            return -1;
        }
        TrieMap<T> next;
        size_t s_len;
        while (is.read(reinterpret_cast<char*>(&s_len), sizeof(s_len))) {
            std::string s_key(s_len, ' ');
            T s_val;
            if (!is.read(s_key.data(), s_len) || !is.read(reinterpret_cast<char*>(&s_val), sizeof(T))) {
                return -1;
            }
            next.add(s_key, s_val);
        }
        std::swap(head, next.head);
        return 0;
    }
};

void normalize(std::string& s) {
    size_t len = s.size();
    for (size_t i = 0; i < len; ++i) {
        if (s[i] >= 'A' && s[i] <= 'Z') {
            s[i] |= 32;
        }
    }
}

int main() {
    std::ios_base::sync_with_stdio(false);
    std::cin.tie(NULL);

    TrieMap<uint64_t> dictionary;
    std::string cmd;

    while (std::cin >> cmd) {
        try {
            if (cmd == "+") {
                std::string k;
                uint64_t v;
                std::cin >> k >> v;
                normalize(k);
                if (dictionary.add(k, v)) {
                    std::cout << "OK\n";
                } else {
                    std::cout << "Exist\n";
                }
            } else if (cmd == "-") {
                std::string k;
                std::cin >> k;
                normalize(k);
                if (dictionary.del_key(k)) {
                    std::cout << "OK\n";
                } else {
                    std::cout << "NoSuchWord\n";
                }
            } else if (cmd == "!") {
                std::string sub, path;
                std::cin >> sub >> path;
                if (sub == "Save") {
                    if (dictionary.store(path) == 0) {
                        std::cout << "OK\n";
                    }
                } else {
                    if (dictionary.restore(path) == 0) {
                        std::cout << "OK\n";
                    }
                }
            } else {
                normalize(cmd);
                uint64_t* res = dictionary.find(cmd);
                if (res) {
                    std::cout << "OK: " << *res << "\n";
                } else {
                    std::cout << "NoSuchWord\n";
                }
            }
        } catch (const std::bad_alloc&) {
            std::cout << "ERROR: Memory failure\n";
        } catch (const std::exception& e) {
            std::cout << "ERROR: " << e.what() << "\n";
        }
    }
    return 0;
}
\end{lstlisting}

\pagebreak

\section{Исследование качества программного обеспечения}

\subsection{Первичное профилирование и обнаружение узких мест}
На начальном этапе разработки было проведено профилирование программы с помощью утилиты \texttt{gprof}. Результаты показали, что значительная часть времени (около 12\% напрямую и до 45\% косвенно через итераторы) тратилась на функцию \texttt{normalize}. Это было вызвано использованием стандартных итераторов строк и функции \texttt{std::tolower}, которая учитывает системные локали, что замедляет обработку каждого символа.

\subsection{Оптимизация кода}
На основе полученных данных были проведены следующие улучшения:
\begin{enumerate}
    \item \textbf{Оптимизация \texttt{normalize}:} Итераторы были заменены на прямой доступ к элементам строки по индексу. Вместо медленной функции \texttt{tolower} был применен побитовый сдвиг (\texttt{s[i] |= 32}), что допустимо для символов латиницы и работает значительно быстрее.
    \item \textbf{Использование \texttt{std::move}:} В основном цикле программы была добавлена семантика перемещения при передаче ключей в структуру данных, что исключило лишние аллокации памяти.
    \item \textbf{Инлайнинг:} Критически важные функции (\texttt{fetch\_bit}) были помечены как \texttt{inline}, что позволило компилятору встроить их код напрямую в места вызова при использовании флага \texttt{-O3}.
\end{enumerate}

\subsection{Повторное профилирование (gprof)}
После внесения изменений было проведено повторное профилирование на тесте из 1,500,000 операций. Результаты Flat Profile:

\begin{center}
\begin{tabular}{|l|c|c|c|c|}
\hline
\% времени & self (сек) & вызовы & мс/вызов & имя функции \\ \hline
55.82 & 2.74 & -- & -- & \texttt{main} (ввод и нормализация) \\ \hline
22.82 & 1.12 & 400,000 & 0.00 & \texttt{TrieMap::add} \\ \hline
20.78 & 1.02 & 400,000 & 0.00 & \texttt{TrieMap::del\_key} \\ \hline
0.41 & 0.02 & 64 & 0.31 & \texttt{frame\_dummy} \\ \hline
\end{tabular}
\end{center}

\textbf{Анализ результатов после оптимизации:}
\begin{enumerate}
    \item Функция \texttt{normalize} и вспомогательные методы работы с битами исчезли из топа профилировщика. Это означает, что они были успешно встроены (\textit{inlined}) компилятором, либо их время выполнения стало пренебрежимо малым.
    \item Основное время (55.8\%) теперь сосредоточено в \texttt{main}, что включает в себя чтение данных из стандартного потока ввода и быструю нормализацию.
    \item Операции добавления и удаления распределены практически поровну (22.8\% и 20.8\%), что подтверждает сбалансированность и предсказуемость дерева PATRICIA.
\end{enumerate}

\subsection{Анализ памяти (AddressSanitizer)}
После оптимизации логики удаления (\texttt{del\_key}) программа была повторно проверена инструментом \texttt{AddressSanitizer} для исключения ошибок сегментации и утечек памяти.
Команда запуска:
\begin{alltt}
g++ -std=c++20 -O3 -g -fsanitize=address main.cpp -o solution
./solution < test.txt
\end{alltt}

\textbf{Результаты:}
Программа успешно обработала 1,500,000 команд без ошибок. Отчет \texttt{LeakSanitizer} показал: \enquote{0 bytes definitely lost}. Все выделенные узлы дерева корректно очищаются в деструкторе через рекурсивный вызов \texttt{deep\_clear}.

\pagebreak

\section{Тест производительности}
Для оценки эффективности реализованной структуры данных было проведено сравнение с \texttt{std::map} из стандартной библиотеки. Тестирование проводилось на наборах данных объемом от 100 до 500 тысяч операций.

\begin{center}
\begin{tabular}{|c|c|c|}
\hline
Кол-во элементов (тыс.) & PATRICIA Tree (мс) & STL \texttt{std::map} (мс) \\ \hline
100 & 150.976 & 364.560 \\ \hline
200 & 331.039 & 889.933 \\ \hline
300 & 744.557 & 1529.501 \\ \hline
400 & 878.791 & 2552.388 \\ \hline
500 & 1744.340 & 3659.964 \\ \hline
\end{tabular}
\end{center}

Ниже представлен график зависимости времени выполнения от количества элементов:

\begin{center}
\begin{tikzpicture}
    \begin{axis}[
        xlabel={Количество элементов (тыс.)},
        ylabel={Время (мс)},
        grid=major,
        legend pos=north west,
        width=12cm, height=8cm,
        title={Сравнение PATRICIA Tree и std::map}
    ]
    
    \addplot[color=blue, mark=square] coordinates {
        (100, 150.97)
        (200, 331.03)
        (300, 744.55)
        (400, 878.79)
        (500, 1744.34)
    };
    \addlegendentry{PATRICIA Tree}

    
    \addplot[color=red, mark=x] coordinates {
        (100, 364.56)
        (200, 889.93)
        (300, 1529.50)
        (400, 2552.38)
        (500, 3659.96)
    };
    \addlegendentry{STL std::map}
    \end{axis}
\end{tikzpicture}
\end{center}

\section{Консоль}
Ниже представлен протокол запуска тестов производительности:
\begin{alltt}
root@5703d59f664e:/workspaces/c_labs/da_lab2# ./solution < test_100k.txt
Patricia Tree: 150.976 ms
STL std::map:  364.560 ms

root@5703d59f664e:/workspaces/c_labs/da_lab2# ./solution < test_200k.txt
Patricia Tree: 331.039 ms
STL std::map:  889.933 ms

root@5703d59f664e:/workspaces/c_labs/da_lab2# ./solution < test_300k.txt
Patricia Tree: 744.557 ms
STL std::map:  1529.501 ms

root@5703d59f664e:/workspaces/c_labs/da_lab2# ./solution < test_400k.txt
Patricia Tree: 878.791 ms
STL std::map:  2552.388 ms

root@5703d59f664e:/workspaces/c_labs/da_lab2# ./solution < test_500k.txt
Patricia Tree: 1744.340 ms
STL std::map:  3659.964 ms
\end{alltt}

\section{Выводы}
В ходе выполнения лабораторной работы была реализована структура данных PATRICIA Tree. Использование инструментов профилирования (\texttt{gprof}) позволило понять, что накладные расходы на работу со стандартными строками \texttt{std::string} сопоставимы с затратами на саму логику дерева. Анализ через \texttt{AddressSanitizer} подтвердил надежность управления памятью.

\pagebreak

\input{../src/literature}

\end{document}